\section{Origin of Revolution}

Revolution is the overthrow of established order, resulting in a state of disorder. Thus, there can be no such thing as a ``conservative revolution'' but rather what is required, as \textbf{Joseph de Maistre} famously claimed, is ``\emph{the opposite of a revolution}``. To understand what Order is in the first place, we turn to \textbf{Dante} in De Monarchia who proposes the logical sequence: Being, Unity, followed by the Good.

\begin{quotex}
Where Being is most, there Unity is greatest; and where Unity is greatest, there Good is also greatest.
\end{quotex}

In the social realm, this Unity is expressed as concord or harmony:

\begin{quotex}
Concord then is the uniform motion of many wills.
\end{quotex}


Firstly, in the individual man himself, there must be harmony and balance of his inner forces. Dante then says the same is true of the family, the city, and the kingdom. This is the state of affairs in a Traditional society where harmony of wills comes from a common mind. The Monarch represents the Will of all in the Kingdom, which then follows in the city and the family. Individual men are expected to harmonize their own states of being, the methods are provided by the spiritual authority through its teachings and the city provides support.

Such a system is characterized by Justice. Now Dante, following Aristotle, points out that

\begin{quotex}
the strongest opponent of Justice is Appetite or \emph{eros}.
\end{quotex}


Hence, Justice reigns only where \emph{eros} is moderated by the Intellect. Hence, the primary tactic of Revolution is to unleash \emph{eros}, on the promise it will make men free. However, just the opposite takes place. Dante explains:

\begin{quotex}
If the Appetite \emph{eros} in any way at all forestalls the Judgment and guides it, then the Judgment cannot be free: it is not its own: it is captive to another power.
\end{quotex}


Thus, the Revolution succeeds by breaking up the harmony of wills, with each man following his own Appetite. This puts them under a power other than the legitimate political and spiritual authorities. In understanding Revolution, it is necessary to identify the agents; this is seldom done. Instead people seem to believe in spontaneous generation, that events arise uncaused out of a rotting carcass. These agents are of three types:

\begin{enumerate}


\item \textbf{Unconscious Agents}. These people absorb the propaganda of Revolution without understanding its larger context.

\item \textbf{Conscious Agents}. These are the powers who deliberating foment revolution. There are two types:

  \begin{enumerate}

  \def\labelenumii{\alph{enumii}.}

  
\item \textbf{Public Agents}. These are publically known through their organizations and platforms.
  
\item \textbf{Secret Agents}. These work behind the scenes, often through a neutral sounding front organization.
  \end{enumerate}

\end{enumerate}


\textbf{Oldstock Ryder} described the art of revolution in his pamphlet The Great Conspiracy.

\begin{quotex}
Revolution is an art but the Revolutionaries would have us believe that it is a natural cataclysm, as inevitable as a volcanic eruption---a spontaneous up-rush of popular revolt against insufferable wrongs... The art of Revolution is that by which a small but well-organized minority compels an unwilling but unorganized majority to submit to the overthrow of the State and the dictatorship of a few professional agitators who grasp power in the name of the People. The method remains the same today as it was in 1789-1793:

\begin{enumerate}


\item First, to create a Revolutionary atmosphere by exploiting the existing grievances or hardships of any part of the population.

\item Secondly, where none exist, popular grievances must be created

\item Thirdly, having thus prepared the stage, demonstrations must be organized which will give the movement an appearance of being a spontaneous uprising of the masses.

\item Fourthly, trade and industry must be hampered and ultimately paralysed by strikes and revolutionary threats, creating widespread unemployment and discontent.

\item Lastly, forces of aliens, criminals, and hooligans must be enlisted and armed to overpower the forces of the State and to terrorize the law-abiding majority into submissions.
\end{enumerate}

And this tragic farce is to be enacted in the name of the whole people, and is applauded as a noble revolt against tyranny and injustice.
\end{quotex}



\flrightit{Posted on 2012-10-16 by Cologero}

\begin{center}* * *\end{center}

\begin{footnotesize}
\begin{sffamily}
\texttt{Avery on 2012-10-16 at 23:20 said:}

These are some good thoughts. Recall that for Evola, ``agents'' of the revolution are not necessarily human, and the forces that sway unconscious agents require careful study. I wish Anders Behring Breivik had read this article.

What does it mean, though, to be the opposite of a revolution? It means being the good guys providing order and authority in the midst of chaos. In the cases of the English, Japanese, and Americans, a certain amount of anti-revolution proved appealing. But the French, Chinese, and Russians proved unable to put eros back in the harness after it had been unleashed; it took decades of tyranny to do that. The mechanics of surviving a revolution are worth looking into.

\hfill

\texttt{Logres on 2012-10-17 at 06:20 said:}

As Dmitri Orlov says, when the tyrants are in power, the ``good people'' cower in their homes, \& will do anything to stay low, right up to the point they are dragged away in the night or stuffed on cattle cars. Solzenhitsyn wondered what would have happened had the secret police known that knocking on someone's door in the night might have shortened their own lifespan. But people just ``went''. Orlov says that the hobbit folk or little people fear ``anarchy'' or furthering chaos, which suggests that they are willing to cling to the Revolution as ``better than nothing''.

This is a fantastic article -- it ties Dante together with ``the ancient city'' and modern times. Take care of your health, sir C.

\hfill

\texttt{fnn on 2012-10-17 at 09:59 said:}

``What does it mean, though, to be the opposite of a revolution? It means being the good guys providing order and authority in the midst of chaos.''

That will get you ten years in prison in the US and maybe a bit less in the UK. But, as far as I can tell, that seems to be what GD is doing in Greece. It doesn't matter that rank-and-file cops are said to support GD, the authorities behind the police clearly don't. And don't forget that leftist Greek mobs throwing Molotov Cocktails

burned to death three bank employees about a year ago.

The first Mayor Daley of Chicago was an unintentional seer on the subject of anarcho-tyranny. He said, ``The policeman isn't there to create disorder; the policeman is there to preserve disorder.''

No endorsement of GD-just an observation.

\hfill

\texttt{Avery Morrow on 2012-10-18 at 08:23 said:}

I agree. This is a key part of understanding GD and it explains why they have such widespread support from police-- the two groups feel the same loyalty to Greece, which more likely than not is centered around the Greek people, and both are endangered by economic and social conditions and desire solidarity.


\hfill

\end{sffamily}
\end{footnotesize}

